\section{Task 3: Combining SPP with Dynamic I-RRIP}
\label{sec:task3}
\subsection{Prefetch-aware replacement code in the supplied archive}
The code archive contains a separate \texttt{replacement/prefetch\_irrip.l2c\_repl} implementation. That file tracks a per-set/per-way \texttt{prefetched} flag, uses the \texttt{PREFETCH} request type, and reads \texttt{CHAMPSIM\_THRESHOLD} at initialization (default 0.25). When a prefetch fills a way, the flag is set and its RRPV is initialized to the maximum value 3, leaving it easy to evict if it is never used. During victim search it first prefers a maximum-RRPV prefetched line; if none qualifies and I-RRIP is active, it next prefers a maximum-RRPV data line, then falls back to any maximum-RRPV line. When a demand hit occurs on a flagged prefetched line, the code counts it as useful, clears the flag, and promotes it using the active RRIP/I-RRIP rules. A prefetch hit itself returns without promoting the line. The policy excludes prefetch requests from the counters used to choose the dynamic mode.

This design follows a reasonable pollution-control rule: speculative lines enter with low protection, while demand reuse gives them normal replacement treatment. The implementation also reports useful and unused-prefetch counts at finalization. These are code-review facts about \texttt{prefetch\_irrip}; they are not, by themselves, evidence of performance improvement.

\subsection{Which Task 3 configuration was measured?}
The Task 3 CSV/JSON metadata and run labels identify the combined
configuration as \texttt{combined-dyn\_irrip}. The supplied runner also
defaults the combined replacement option to \texttt{dyn\_irrip}. In the
provided source archive, \texttt{dyn\_irrip.l2c\_repl} does not implement
the per-way prefetch tracking/eviction logic described above. Consequently,
the measurements below are labeled \textbf{SPP + dyn\_irrip (as recorded)};
they cannot be presented as measured results for the separate prefetch-aware
\texttt{prefetch\_irrip} file. The Task 3 requirement should be re-evaluated
after the prefetch-aware logic is integrated into the specified dynamic policy
and the runs are repeated.

\subsection{Measured results: IPC, speedup, and L1D behavior}
\begin{table}[H]\centering\scriptsize
\caption{Task 3 data as recorded. Speedup is relative to no-prefetch + LRU for the same trace.}
\resizebox{\linewidth}{!}{\begin{tabular}{llrrrrr}
\toprule Trace & Configuration & IPC & Speedup & Accuracy (\%) & Coverage (\%) & L1D load MPKI \\
\midrule
Tr1.gz & Baseline (LRU, no prefetch) & 0.232345 & 1.000000 & N/A & N/A & 10.6780 \\
Tr1.gz & SPP only & 0.232577 & 1.000999 & 41.33 & 1.23 & 10.5465 \\
Tr1.gz & Dynamic I-RRIP only & 0.232281 & 0.999725 & N/A & N/A & 10.6788 \\
Tr1.gz & SPP + \texttt{dyn\_irrip} (as recorded) & 0.232482 & 1.000590 & 41.07 & 1.22 & 10.5476 \\
Tr2.gz & Baseline (LRU, no prefetch) & 0.251158 & 1.000000 & N/A & N/A & 13.5379 \\
Tr2.gz & SPP only & 0.251094 & 0.999745 & 33.54 & 0.27 & 13.5014 \\
Tr2.gz & Dynamic I-RRIP only & 0.256285 & 1.020413 & N/A & N/A & 13.5412 \\
Tr2.gz & SPP + \texttt{dyn\_irrip} (as recorded) & 0.256303 & 1.020485 & 33.34 & 0.23 & 13.5064 \\
Tr3.gz & Baseline (LRU, no prefetch) & 0.382065 & 1.000000 & N/A & N/A & 3.4034 \\
Tr3.gz & SPP only & 0.382139 & 1.000194 & 8.51 & 0.28 & 3.3938 \\
Tr3.gz & Dynamic I-RRIP only & 0.381742 & 0.999155 & N/A & N/A & 3.4033 \\
Tr3.gz & SPP + \texttt{dyn\_irrip} (as recorded) & 0.381886 & 0.999531 & 8.52 & 0.29 & 3.3937 \\
\bottomrule\end{tabular}}
\end{table}

\subsection{L2 behavior}
\begin{table}[H]\centering\scriptsize
\caption{L2 statistics by trace/configuration. These use the separate aggregate L2C and instruction/data counters described in Section 1.}
\resizebox{\linewidth}{!}{\begin{tabular}{llrrr}
\toprule Trace & Configuration & L2C load MPKI & L2 instruction-load MPKI & L2 data-load MPKI \\
\midrule
Tr1.gz & Baseline (LRU, no prefetch) & 2.2774 & 10.3377 & 7.3580 \\
Tr1.gz & SPP only & 2.2805 & 10.3556 & 7.2210 \\
Tr1.gz & Dynamic I-RRIP only & 2.3941 & 9.9094 & 7.6956 \\
Tr1.gz & SPP + \texttt{dyn\_irrip} (as recorded) & 2.3894 & 9.9151 & 7.5314 \\
Tr2.gz & Baseline (LRU, no prefetch) & 2.2437 & 18.8533 & 8.5035 \\
Tr2.gz & SPP only & 2.2457 & 18.9067 & 8.4938 \\
Tr2.gz & Dynamic I-RRIP only & 2.3631 & 13.3453 & 9.1221 \\
Tr2.gz & SPP + \texttt{dyn\_irrip} (as recorded) & 2.3648 & 13.3179 & 9.1141 \\
Tr3.gz & Baseline (LRU, no prefetch) & 0.6972 & 1.1241 & 2.8302 \\
Tr3.gz & SPP only & 0.7033 & 1.1386 & 2.8119 \\
Tr3.gz & Dynamic I-RRIP only & 0.7235 & 1.1464 & 2.8764 \\
Tr3.gz & SPP + \texttt{dyn\_irrip} (as recorded) & 0.7275 & 1.1494 & 2.8502 \\
\bottomrule\end{tabular}}
\end{table}

\plotpair{plots/task3/ipc.png}{IPC}{plots/task3/speedup.png}{Speedup}{Task 3 measured configuration comparison}{fig:t3-ipc-speedup}
\plotpair{plots/task3/l1d_load_mpki.png}{L1D load MPKI}{plots/task3/l1d_coverage_pct.png}{L1D coverage}{Task 3 L1D miss and coverage behavior}{fig:t3-l1d}
\plotpair{plots/task3/prefetch_accuracy_pct.png}{Prefetch accuracy}{plots/task3/l2c_load_mpki.png}{L2C load MPKI}{Task 3 prefetch accuracy and aggregate L2C metric}{fig:t3-accuracy-l2c}
\plotpair{plots/task3/l2i_load_mpki.png}{L2 instruction-load MPKI}{plots/task3/l2d_load_mpki.png}{L2 data-load MPKI}{Task 3 L2 instruction and data behavior}{fig:t3-l2i-l2d}
\plotpair{plots/task3/ipc_change_vs_baseline_pct.png}{IPC change vs. baseline (\%)}{plots/task3/ipc_gain_vs_best_standalone_pct.png}{IPC gain vs. best standalone (\%)}{Additional Task 3 relative-performance views}{fig:t3-relative}

\subsection{Does the measured combined configuration meet the IPC goal?}
The answer is trace-dependent even for the configuration recorded in the CSVs. On Tr1, the combined IPC is 0.232482, above the baseline and dynamic-only result but below SPP-only (0.232577). On Tr2, combined IPC is 0.256303, just above dynamic-only (0.256285) and clearly above SPP-only (0.251094); the difference over dynamic-only is only about 0.007\%. On Tr3, combined IPC is 0.381886, above dynamic-only (0.381742) but below SPP-only (0.382139) and below the no-prefetch baseline (0.382065). Therefore the recorded configuration does not outperform both standalone methods on all traces.

The results also show why no single metric is enough. Prefetch accuracy is approximately 41.07\%, 33.34\%, and 8.52\% for the combined configuration on Tr1--Tr3, while coverage is about 1.22\%, 0.23\%, and 0.29\%. L1D MPKI is modestly lower than baseline on all traces, but aggregate L2C MPKI is higher on all three. At the same time, instruction-load MPKI is considerably lower on Tr2, while data MPKI is higher than the no-prefetch baseline on all traces. Because the combined runs used \texttt{dyn\_irrip}, these observations are about the measured SPP + dynamic replacement configuration and must not be used to infer the effect of prefetch-aware victim selection.
